% 第 15 题（2）解析：地面系中小球的轨迹（TikZ 重绘）
% 实线为小球的轨迹椭圆（中心在 ma/(m0+m) 处），虚线为某时刻凹槽所在的椭圆；
% 两椭圆交于小球所在位置。标注 x、y、x1'、x2'
\begin{tikzpicture}[font=\small, >=Stealth]
  % 坐标轴
  \draw[->, thick] (-1.45,0) -- (4.7,0) node[below]{$x$};
  \draw[->, thick] (0,-1.55) -- (0,1.05) node[left]{$y$};
  \node[below left] at (0,0) {$O$};
  % 小球的轨迹椭圆（实线，中心 (1,0)，半轴 2 与 1.2）
  \draw[thick] (3,0) arc[start angle=0, end angle=-180, x radius=2, y radius=1.2];
  % 凹槽此刻所在的椭圆（虚线，中心 (0.6,0)，半轴 3 与 1.2）
  \draw[dashed] (3.6,0) arc[start angle=0, end angle=-180, x radius=3, y radius=1.2];
  % 运动坐标系
  \draw[dashed] (0.6,-1.4) -- (0.6,0.95);
  \node[above right] at (0.6,0.95) {$y'$};
  % 小球
  \fill[gray!40] (1.8,-1.1) circle (0.12);
  \draw[thick] (1.8,-1.1) circle (0.12);
  \fill (3,0) circle (1.2pt);
  % 位移标注
  \draw[<->] (0,0.32) -- (0.6,0.32) node[midway, above]{$x_1'$};
  \draw[<->] (1.8,0.32) -- (3,0.32) node[midway, above]{$x_2'$};
  \draw[<->] (0,-1.45) -- (1.8,-1.45) node[midway, below]{$x$};
  \draw[dashed] (1.8,-1.1) -- (1.8,0);
  \node[right] at (1.84,-0.62) {$y$};
\end{tikzpicture}
